% !TEX root = ../main.tex

\chapter{总结与展望}
\label{chap:conclusion}

\section{工作总结}

本文面向 6G 空天地一体化网络中的大规模 LEO 卫星星座，研究分布式网络管控下的端到端用户 QoS 保障问题。从全链路业务承载过程看，端到端 QoS 保障可分解为感知、接入、资源分配与传输四个环节；从系统能力指标看，其核心在于同时支撑通信系统可靠性、资源分配效率与接入公平性。在缺乏可靠全局状态与中心化协调的条件下，上述目标具体表现为四类难题：退化遥测下节点难以构建可信的本地多跳网络视图；无中心协调的海量并发切换使个体切换效率与系统级长期公平相互冲突；多用户竞争与业务突发使资源决策对用户级 QoS 的影响具有滞后性和随机性，受限的物理资源易被当前预算下难以满足 QoS 目标的业务流消耗；星上计算资源受限使语义传输难以简单依靠提升发射端编码模型的表达能力，或增强接收端恢复模型的纠错与恢复能力来改善任务性能。

针对上述难题，本文的核心思路是以本地状态校准、有限邻域交互与跨节点协作替代全局状态同步与集中式求解，沿端到端路径依次设计四项可在星上在线执行的关键技术：在感知环节，构建多源置信度调控的分布式局部网络状态感知框架，为分布式路由提供可信的本地网络视图；在接入环节，构建基于局部公平性估计的分布式协作接入切换机制，协调个体切换效率与系统长期公平；在资源分配环节，构建基于生成式网络态势感知的切片--MAC 分层协作优化框架，协调资源使用效率与用户级 QoS 保障；在传输环节，构建基于任务反馈与协同增量重传的轻量化语义传输机制，协调端到端传输性能与星上硬件开销。四项技术覆盖端到端传输链路的各个环节，共同构成从状态获取、接入建立、资源分配到业务传输的闭环管控体系。本文的主要工作与所取得的成果归纳如下。

\textbf{（1）多源置信度调控的分布式局部网络状态感知框架。} 在端到端传输链路的感知环节，退化遥测会污染状态聚合，削弱路由决策的可靠性。本文提出置信度自适应局部分布式处理（CA-LDP）框架，建立退化遥测下先校准、后聚合的可信本地状态感知机制。第一项成果是置信度自适应的图信号修复机制，在时间相关性、空间相关性与信息新鲜度等多置信源之间通过归因分析校准陈旧与缺失遥测，实现精确的网络状态恢复；第二项成果是受置信度调制的图信号聚合机制，抑制多跳传播中低置信节点对可靠节点的污染。在由 1584 颗卫星构成的 Starlink 真实星座、九类受控遥测退化场景下，CA-LDP 相较代表性基线将平均传输速率提升 $4.8\%$--$10.8\%$，端到端时延降低 $6.7\%$--$11.1\%$，并维持接近 $100\%$ 的路由成功率。上述结果表明，CA-LDP 能够在遥测陈旧、缺失与噪声并存时维持可信的本地多跳网络视图，解决退化遥测条件下分布式路由缺少可靠状态输入的问题。

\textbf{（2）基于局部公平性估计的分布式协作接入切换机制。} 在接入环节，无中心协调的海量并发切换使个体接入容量与系统级长期公平相互冲突。本文提出公平性引导的多智能体近端策略优化（FG-MAPPO）方法，在无全局协调条件下实现公平感知的分布式协作切换。第一项成果是基于广义 Gini 社会福利函数的优势修正机制，将系统级公平目标转化为用户可在邻域内本地估计的优势修正量，实现低开销的邻域公平性估计；第二项成果是吞吐与公平分立的双目标学习结构，在缓解多智能体非平稳性的同时保留个体速率性能。在仅交互邻域信息的条件下，FG-MAPPO 的广义 Gini 公平性指标较具备全局信息的集中式匈牙利算法提升 $3.48\%$、平均速率提升 $6.91\%$；相较分布式强化学习基线，其切换成功率与负载均衡性能最高分别提升 $15.65\%$ 与 $30.91\%$。上述结果表明，FG-MAPPO 能够在缺乏集中式协调的海量并发切换中同时维持接入效率、服务连续性和长期公平，解决个体接入容量与系统级接入公平难以兼顾的问题。

\textbf{（3）基于生成式网络态势感知的切片--MAC 分层协作优化框架。} 在资源分配环节，多用户竞争与业务突发使资源决策对用户级 QoS 的影响滞后且随机波动，而反应式资源分配通常滞后于队列积压和 QoS 违背，使切片层与 MAC 层难以形成面向未来业务状态的跨层协同，无法实现用户级 QoS 的长期保障。本文提出双时间尺度预测调度（DTPS）框架，以带宽与功率预算为跨层接口、以 QoS 可行性为协同准则，将长期 QoS 保障分解为切片层的 QoS 可行域保障与 MAC 层的实时调度。第一项成果是基于世界模型的生成式切片资源预留机制，其生成式态势感知包含训练和在线推理两个要点：在训练阶段，通过预测轨迹与真实观测轨迹对齐，使模型生成的业务到达、队列演化和信道变化尽可能接近真实无线网络状态，从而学习流量和信道变化的物理统计规律；在在线推理阶段，模型利用实时观测更新内部网络状态表示，并在线调整对未来无线网络状态和用户级 QoS 演化的预测，进而提前确定切片带宽与功率预算。第二项成果是 QCI 优先级背包调度机制，在预算约束的可行域内估计各候选流满足速率--时延联合门限所需的最小 RBG--功率开销，将受限资源优先配置给 QoS 目标可满足的业务流。在 Starlink 规模星座仿真中，DTPS 在过载业务下实现 $80.8\%$ 的用户级 QoS 满足率，显著高于 TDRL-RSUA、RL-JRSS、LACO 与 DQN-MAC 等基线的 $33.0\%$--$56.3\%$，并将平均时延保持在 $78.3$~ms；突发不确定性增强时，其 QoS 满足率始终维持在 $97.6\%$ 以上。上述结果表明，DTPS 能够在无线资源受限且业务快速波动的动态 LEO 星座中维持用户级 QoS 满足能力，解决用户级 QoS 难以长期保障的问题。

\textbf{（4）基于任务反馈与协同增量重传的轻量化语义传输机制。} 在传输环节，星上计算资源受限使语义传输难以简单依靠提升发射端编码模型的表达能力，或增强接收端恢复模型的纠错与恢复能力来改善任务性能。本文提出协同增量语义混合自动重传请求（CIS-HARQ）机制，将每一轮 HARQ 由完整语义重传重构为任务反馈引导的选择性增量编码与跨轮合并，确立以扩展通信轮次和反馈协作替代扩展星上计算的可靠传输路径。第一项成果是基于任务反馈的协同增量重传框架，建立任务感知接收端与轻量发射端之间的闭环协作：接收端根据当前接收语义状态与任务所需判别特征之间的差距生成反馈，发射端依据反馈调整后续增量编码，接收端再将新接收增量与历史语义状态合并，逐轮恢复任务相关特征。第二项成果是该框架的低开销实现方法，包括基于任务判决裕量归因的语义分块价值排序、基于熵估计的增量码率约束、发射端基于 SFT 的反馈条件化编码，以及接收端掩码引导的门控语义合并。在 ImageNet-100、CIFAR-100 与 VeRi-776 数据集、Rician 衰落信道、信噪比 $5$~dB 与信道带宽比约 $0.02$ 的紧带宽条件下，CIS-HARQ 经三轮增量传输将下游任务精度由同带宽基线的 $35.7\%$ 提升至 $65.0\%$；其以约 $2.2$~M 参数与毫秒级星上推理时延，较生成式扩散接收端减少约 $2500$ 倍参数量，任务精度差距控制在约 $7\%$ 以内，表明轻量化增量重传能够在受限星上算力下逼近高复杂度语义模型的任务性能，并在不同数据分布与多任务接收端上保持一致的增量传输增益。

综上，本文构建了覆盖可信感知、公平接入、高效资源分配与可靠传输的端到端闭环管控体系，在星座规模仿真与语义传输实验中，分别缓解了退化遥测下本地状态不可信、无中心切换下接入公平性不足、业务突发下长期用户级 QoS 难以维持、紧带宽下轻量语义传输精度不足等问题。在方法论层面，本文将大规模 LEO 星座资源管控由分层独立优化推进至面向端到端 QoS 的分布式协同管控，论证了在缺乏可靠全局状态与中心化协调的条件下，依靠本地状态校准与受限邻域协作仍可逼近乃至超越集中式方案的性能；在工程层面，所提机制均以固定计算预算在受限星上载荷上在线执行，并与现行 3GPP 协议流程兼容，为大规模 LEO 星座在分布式网络管控下实现端到端用户 QoS 保障提供了工程可行的技术途径。

\section{研究展望}

本文围绕单层 LEO 星座中的状态感知、接入切换、无线资源分配与语义可靠传输，构建了面向端到端 QoS 保障的分布式协同管控方法。面向未来大规模 LEO 星座向多层异构、跨系统协同、渐进部署和多样化业务演进，相关方法仍需在网络态势表达、网络功能融合、管控能力演进和语义可靠传输等层面继续拓展。

在网络状态感知方面，本文主要针对退化遥测条件下的本地多跳状态恢复与可信聚合，解决了单层星座内分布式路由缺少可靠状态输入的问题。未来面向 LEO 与中高轨星座、高空平台及地面网络融合的多层异构组网，状态感知需要由局部邻域视图进一步扩展为跨层、跨域和跨运营主体的网络态势表达。时空动态图建模、网络数字孪生与网络状态语义化压缩可作为重要方向，用于研究多尺度网络状态的表征、预测与压缩问题；同时，跨系统状态感知还需要考虑可信信息融合与安全交互机制，使不同运营主体和不同制式网络能够在有限信息共享条件下形成一致、可验证的网络状态认知。

在网络功能部署方面，本文主要围绕接入网侧的状态感知、接入切换、无线资源分配与业务传输展开，尚未进一步处理核心网功能与接入网功能在星上节点中的联合部署问题。未来再生式 LEO 卫星网络需要将核心网控制面、用户面、接入网协议栈和星间回传功能纳入统一分析框架，研究网络功能下沉、控制面分布和接入--回传资源耦合之间的内在关系。该方向的核心在于刻画网络功能融合对端到端时延、业务连续性、控制可靠性和星上资源消耗的影响，为星上网络功能放置与资源分配提供理论依据。

在资源管控演进方面，本文所提方法在给定星座构型和业务场景下验证了学习型分布式管控对接入公平、资源分配和语义传输的增益。面向分阶段建设和长期运行的大规模星座，网络规模、覆盖结构和业务分布会随时间持续变化。因此，未来的学习型管控机制需要具备持续演进和增量学习能力，在新卫星、新覆盖区域和新业务类型接入时实现快速适配，同时维持既有区域的服务质量稳定。围绕这一目标，模型迁移、在线学习、泛化性能和稳定性--可塑性权衡等基础问题仍有待进一步研究。

在语义通信方面，本文通过任务反馈与协同增量重传验证了轻量化语义 HARQ 在紧带宽和受限星上计算资源下逼近高复杂度语义模型任务性能的可行性。未来需要进一步从通信理论层面分析面向任务的语义可靠传输，建立语义失真、语义错误概率和反馈信息价值等度量，并刻画语义 HARQ 在可靠性、时延和资源开销方面的性能界。同时，语义传输也应由单次视觉任务扩展至多模态业务，并进一步利用跨任务、跨时间和跨节点的语义相关性开展协同传输优化，从而提升受限链路下任务相关信息的恢复效率。
